\section{Statement of results}\label{section2}

Let $a_j$ be distinct complex numbers, $a_j\not\in\{0,1\}$ for $1\leq j\leq k$. Let $m_j$ be positive integers, and $\alpha$, $\beta$, $\gamma$, $\delta$ complex numbers. Consider the differential equation
\begin{gather}\label{1a}
y''-\left( \frac\alpha z+\frac\beta{z-1}+\sum_{i=1}^k \frac{m_i}{z-a_i}\right) y'+\frac{\gamma\delta z^k+\cdots{}}{z(z-1)\prod\limits_{i=1}^k(z-a_i)}y = 0,
\end{gather}
where $\gamma+\delta=\alpha+\beta+1+\sum\limits_{j=1}^km_j$ (Fuchs's condition). The Riemann scheme of this equation is
\begin{gather*}
\left(\begin{matrix}
0&1&a_1&\ldots&a_k&\infty\\
0&0&0&\ldots&0&-\gamma\\
\alpha+1&\beta+1&m_1+1&\ldots&m_k+1&-\delta
\end{matrix}\right).
\end{gather*}
We recall that the first line of the Riemann scheme contains the singularities, and the other two the exponents at each singularity~\cite{Mellin}. Set $d=\sum\limits_{j=1}^km_j$. Assume that the singularities at \smash{$a_1\lc a_k$} are all apparent and
\begin{gather}\label{0d}
\gamma, \delta, \gamma-\alpha-1, \delta-\alpha-1 \not\in \{0,1\lc d-1\}.
\end{gather}

\begin{thm}\label{theorem1}
Under these conditions, there exists a polynomial $Q$ of degree $d=\sum\limits_{j=1}^km_j$ such that every solution of \eqref{1a} has the form $Q\big(z\frac{{\rm d}}{{\rm d}z}\big)F(z) $, where $F$ is a solution of the hypergeometric equation
\begin{gather}\label{2a}
F''-\left( \frac\alpha z+\frac{\beta+d}{z-1} \right) F'+\frac{\gamma\delta}{z(z-1)} F = 0.
\end{gather}
In particular, the monodromy group of equation \eqref{1a} coincides with that of the hypergeometric equation~\eqref{2a}.
\end{thm}

\begin{Remark}\label{remark1} Clearly, conditions \eqref{0d} imply that for any non-zero polynomial $p$ of degree at most $d-1$, the functions $p(z)$ and $z^{\alpha+1}p(z)$ are not solutions of the hypergeometric equation~\eqref{2a}.
\end{Remark}

\begin{Remark}\label{remark2} Under conditions \eqref{0d}, for any non-zero polynomial $Q$ of degree at most $d$ and any non-zero solution $F$ of \eqref{2a}, the function $Q\big(z\frac{{\rm d}}{{\rm d}z}\big)F(z)$ is non-zero. Indeed, a solution $F$ of~\eqref{2a} such that $Q\big(z\frac{{\rm d}}{{\rm d}z}\big)F(z)\equiv 0$, has the following form: there are polynomials $p_1$, $p_2$ of degrees at most $d-1$ and constants $A,B$, such that
\begin{alignat*}{3}
& F(z) = p_1(z) + z^\alpha p_2(z)\qquad && \hbox{\rm if $\alpha$ is not an integer},&\\
& F(z) = p_1(z) + z^\alpha p_2(z) (A+B\log z)\qquad && \hbox{\rm if $\alpha$ is a non-negative integer},&\\
& F(z) = z^\alpha p_1(z) + p_2(z) (A+B\log z)\qquad && \hbox{\rm if $\alpha$ is a negative integer}. &
\end{alignat*}
Then in the first case, both $p_1(z)$ and $z^\alpha p_2(z)$ are solutions of equation~\eqref{2a} in contradiction with Remark~\ref{remark1}. In the second case, $z^\alpha p_2(z)$ is a solution of \eqref{2a} if $B\ne 0$, while for $B=0$, $F(z)$ is a polynomial and $ \deg F<d$ since $Q\big(z\frac{{\rm d}}{{\rm d}z}\big)F(z)=0$. Both options contradict Remark~\ref{remark1}. The third case can be worked out similarly.
\end{Remark}

\begin{thm}\label{theorem2}
For every $\alpha$, $\beta$, $\gamma$, $\delta$, $m_1\lc m_k$, $a_1\lc a_k$, there exist at least one and at most $\prod\limits_{j=1}^k(m_j+1)$ equations~\eqref{1a} with apparent singularities at all points $a_j$. For generic $\alpha$, $\beta$, $\gamma$, $\delta$, $a_1\lc a_k$, the equality holds.
\end{thm}
